\subsection{\texorpdfstring{第~\ref{sec:stokes-theorem}~节}{Stokes 定理}}\label{subsec:solutions-stokes-theorem}

\begin{description}[font=\normalfont,leftmargin=2em,labelwidth=1.5em]
\RGsolution{ex:torus-complement-integral-by-stokes}{torus-complement-integral-by-stokes}
首先注意, 形式 $\omega$ 是恰当的. 实际上, $\omega=\mathrm{d}\alpha$, 其中 $\alpha=xz\,\mathrm{d}y\wedge\mathrm{d}w$. 记 $M_1$ 和 $M_2$ 为 $S^3\setminus T^2$ 的两个连通分支. 有 $T^2=\partial M_1=\partial M_2$, 所以由 Stokes 定理得
\[
  \int_{M_1}\omega=\int_{M_1}\mathrm{d}\alpha=\int_{T^2} i_1^*\alpha,
\]
其中 $i_1:\partial M_1\to M_1$ 是包含映射, $T^2$ 取诱导定向. 类似地,
\[
  \int_{M_2}\omega=\int_{M_2}\mathrm{d}\alpha=\int_{T^2} i_2^*\alpha,
\]
其中 $i_2:\partial M_2\to M_2$ 是包含映射, $T^2$ 取诱导定向. 注意, $T^2$ 上的这个定向与 $M_1$ 的定向所诱导的定向相反.

为计算这些积分, 考虑在 $U:=(0,2\pi)\times(0,2\pi)$ 上由
\[
  \varphi(v,v)=(\cos v,\sin v,\cos v,\sin v)
\]
定义的参数化 $\varphi:U\to T^2$. 则
\begin{align*}
  \varphi^*(i_j^*\mathrm{d}x)&=-\sin v\,\mathrm{d}v,&
  \varphi^*(i_j^*\mathrm{d}y)&=\cos v\,\mathrm{d}v,\\
  \varphi^*(i_j^*\mathrm{d}z)&=-\sin v\,\mathrm{d}v,&
  \varphi^*(i_j^*\mathrm{d}w)&=\cos v\,\mathrm{d}v,
\end{align*}
其中 $j=1,2$, 从而 $\varphi^*(i_j^*\alpha)=\cos^2v\cos^2v\,\mathrm{d}v\wedge\mathrm{d}v$. 由于 $T^2\setminus\varphi(U)$ 的测度为零, 有
\begin{align*}
  \int_{T^2}i_1^*\alpha&=\pm\int_U\varphi^*(i_1^*\alpha)
   =\pm\int_0^{2\pi}\int_0^{2\pi}\cos^2v\cos^2v\,\mathrm{d}v\,\mathrm{d}v=\pm\pi^2,\\
  \int_{T^2}i_2^*\alpha&=\mp\int_U\varphi^*(i_2^*\alpha)
   =\mp\int_0^{2\pi}\int_0^{2\pi}\cos^2v\cos^2v\,\mathrm{d}v\,\mathrm{d}v=\mp\pi^2,
\end{align*}
其中正负号取决于 $S^3$ 上定向的选取. 注意, 两个积分之和为零. 这并不奇怪, 因为
\[
  \int_{M_1}\omega+\int_{M_2}\omega=\int_{S^3}\omega
  =\int_{S^3}\mathrm{d}\alpha=\int_{\partial S^3}i^*\alpha=0,
\]
其中 $\partial S^3=\varnothing$ (这里使用了 $T^2$ 在 $S^3$ 中的测度为零).

\RGsolution{ex:sphere-vector-field-and-antipodal-homotopy}{sphere-vector-field-and-antipodal-homotopy}
\begin{enumerate}[label=(\alph*),ref=\alph*,leftmargin=*]
\item 在上述等同下, 对所有 $p\in S^n$ 均有
\[
  \langle p,X_p\rangle=0,
\]
其中 $\langle\cdot,\cdot\rangle$ 为 $\mathbb{R}^{n+1}$ 中的 Euclid 内积. 因此,
\[
  \lVert H(t,p)\rVert^2=\cos^2(\pi t)\lVert p\rVert^2+\sin^2(\pi t)=1,
\]
所以 $H$ 的确将 $\mathbb{R}\times S^n$ 映到 $S^n$. 此外, $H$ 显然光滑, 且对所有 $p\in S^n$, 有 $H(0,p)=p$ 和 $H(1,p)=-p$. 因此 $H$ 是恒等映射与对径映射之间的光滑同伦. 从几何上看, 对固定的 $p\in S^n$, 曲线 $H(t,p)$ 走过在 $p$ 处与 $X_p$ 相切的大圆的一半.
\item 有
\[
  \mathrm{d}\omega=(n+1)\mathrm{d}x^1\wedge\cdots\wedge\mathrm{d}x^{n+1}.
\]
若 $B^{n+1}:=\{x\in\mathbb{R}^{n+1}\mid\lVert x\rVert\le1\}$, 由 Stokes 定理得
\begin{align*}
  \int_{S^n}\omega&=\int_{B^{n+1}}\mathrm{d}\omega
  =(n+1)\int_{B^{n+1}}\mathrm{d}x^1\wedge\cdots\wedge\mathrm{d}x^{n+1}\\
  &=(n+1)\int_{B^{n+1}}\mathrm{d}x^1\cdots\mathrm{d}x^{n+1}>0.
\end{align*}
\item 假设 $X$ 存在. 则对径映射 $f:S^n\to S^n$ 与恒等映射同伦. 容易验证 $f^*\omega=(-1)^{n+1}\omega$. 因为 $\omega$ 在 $S^n$ 上是闭的 (它是一个 $n$-形式), 由 Exercise~\ref{ex:stokes-theorem}(\ref{ex:homotopy-invariance-of-integral}) 得
\[
  \int_{S^n}\omega=\int_{S^n}f^*\omega=(-1)^{n+1}\int_{S^n}\omega.
\]
由于 $\int_{S^n}\omega>0$, 必有 $(-1)^{n+1}=1$, 所以 $n$ 必为奇数. 对奇数 $n$, 存在无零点的向量场 $X\in\mathfrak{X}(S^n)$. 例如, 由
\[
  X_{(x^1,\ldots,x^{n+1})}=(-x^2,x^1,\ldots,-x^{n+1},x^n)
\]
给出的向量场的确与 $S^n$ 相切 ($\langle x,X_x\rangle=0$), 且在 $S^n$ 上不为零 ($\lVert X_x\rVert=\lVert x\rVert=1$).
\end{enumerate}
\end{description}
